% Generated draft: requires mathematical and source review.
\documentclass[11pt]{article}
\usepackage[margin=25mm]{geometry}
\usepackage{fontspec}
\setmainfont{Latin Modern Roman}
\usepackage{amsmath,amssymb,mathtools,unicode-math}
\setmathfont{Latin Modern Math}
\usepackage{enumitem,needspace,xurl,hyperref}
\setlength{\parindent}{0pt}
\setlength{\parskip}{5pt}
\setlength{\emergencystretch}{3em}
\newcommand{\N}{\mathbb{N}}
\newcommand{\Z}{\mathbb{Z}}
\newcommand{\Q}{\mathbb{Q}}
\newcommand{\R}{\mathbb{R}}
\newcommand{\C}{\mathbb{C}}
\newcommand{\F}{\mathbb{F}}
\newcommand{\A}{\mathbb{A}}
\newcommand{\PP}{\mathbb{P}}
\newcommand{\E}{\mathbb{E}}
\newcommand{\eps}{\varepsilon}
\newcommand{\dd}{\,\mathrm{d}}
\begin{document}
\section*{UnsolvedMath ID: 30002399}
\textbf{Source problem code:} catalogue entry
\subsection{Definitions and mathematical statement}
The supplied record asks two questions about categorifications of link invariants. It does not specify the admissible categorifications, their equivalences, or the conventions for the link invariants. Consequently, the following is a conditional formulation schema, rather than a fully specified conjecture.

Fix a class \(\mathcal L\) of links modulo the intended link equivalence, a coefficient ring \(R\), and a link invariant \(I\colon\mathcal L\to R\). The orientation, framing, and normalization conventions are unresolved in the supplied record. Once a collection \(\mathcal A\) of admissibility requirements and a decategorification procedure \(\operatorname{Dec}\) have been specified, write
\[
\mathfrak C_{\mathcal A}(I)
 =\{\mathcal H: \mathcal H\text{ satisfies }\mathcal A
       \text{ and }\operatorname{Dec}(\mathcal H)=I\}.
\]
Here \(\mathcal H\) denotes an entire categorification theory for the chosen links, and the equality means equality of the resulting link invariants. Let \(\simeq\) denote the intended equivalence of such theories, which also remains unspecified.

\begin{enumerate}
\item For a given link invariant \(I\), under what admissibility requirements and notion of equivalence, if any, does uniqueness hold in the sense that
\[
\forall\mathcal H_1,\mathcal H_2\in\mathfrak C_{\mathcal A}(I),
\qquad \mathcal H_1\simeq\mathcal H_2?
\]
This formula expresses uniqueness conditional on existence. The source does not say whether existence for the given invariant is assumed, nor whether the question concerns every link invariant or a particular intended family.
\item Does the Kauffman polynomial admit a categorification? If \(K\) denotes the Kauffman polynomial with the intended conventions, the corresponding existence question is
\[
\mathfrak C_{\mathcal A}(K)\ne\varnothing.
\]
The polynomial variant, normalization, coefficient ring, grading conventions, decategorification procedure, and admissibility requirements needed to interpret this assertion are not specified.
\end{enumerate}

The source wording that categorifications might ``become unique'' does not identify any additional hypotheses or operation under which uniqueness is intended. Neither the parameters above nor a relationship between the two questions can therefore be fixed from the supplied record alone. The supplied dated curation labels the problem open and distinguishes uniqueness under restricted tensor-product categorification axioms from the broader question; that label and assessment are not independently verified here.
\subsection{Short English statement}
The questions ask whether categorifications of a fixed link invariant are unique up to an appropriate equivalence, and whether the Kauffman polynomial has a categorification. A precise conjecture requires additional choices about admissible theories, equivalence, decategorification, and polynomial conventions. The supplied open-status assessment is dated curation, not an independently verified mathematical conclusion.
\subsection{Sources}
\begin{itemize}
\item Local UnsolvedMath registry: permanent ID \(30002399\), title \emph{Uniqueness and Existence of Link-Invariant Categorifications}, source problem code \texttt{catalogue entry}, in the Oberwolfach Reports Open Problems collection. No individual proposer is identified; the proposed year is recorded as 2013.
\item Supplied source citation: \emph{Matrix Factorizations in Algebra, Geometry, and Physics} (2014), Oberwolfach Reports, pp.\ 2501--2552, DOI \texttt{10.4171/owr/2013/44}. \url{https://doi.org/10.4171/owr/2013/44}. The supplied background attributes the question to workshop open-problem discussions documented in Oberwolfach Reports.
\item Supplied literature source, labeled ``Oberwolfach Report'': \url{https://ems.press/content/serial-article-files/46474}.
\item Supplied literature source, labeled ``Losev--Webster'': \url{https://arxiv.org/abs/1303.1336}.
\item Literature assessment dated exactly \texttt{2026-08-21}: the supplied curation labels the problem open and states that uniqueness is proved under restricted tensor-product categorification axioms, while the broader uniqueness question for the stated link invariant, including the indicated case, is not settled. This is supplied dated curation, not independently verified; the record itself requests independent verification of time-sensitive status and cited claims.
\item The registry marks the statement as exact and reports that its mathematical content was checked against the official Oberwolfach report, with notation, context, and question form normalized. These are supplied verification claims; no source was independently consulted for this draft.
\end{itemize}
\end{document}
